% ============================================================
%  CapMismatch: Measuring Capability-Description Inconsistency
%  and Validating Source-Sink Attack Risks in LLM Agent
%  Skill Marketplaces
%
%  Target: NDSS 2027
%  Format: NDSS official template (IEEE two-column)
%  Body:   ≤ 10 pages (excl. references & appendices per NDSS rules)
% ============================================================

\documentclass[letterpaper,twocolumn,10pt]{article}

% ── NDSS 2027 template packages ──────────────────────────────────────────────
\usepackage{ndss2025}        % official NDSS style file (rename if needed)
\usepackage[utf8]{inputenc}
\usepackage[T1]{fontenc}
\usepackage{microtype}
\usepackage{times}

% ── Standard packages ────────────────────────────────────────────────────────
\usepackage{graphicx}
\usepackage{amsmath,amssymb,amsthm}
\usepackage{booktabs}
\usepackage{multirow}
\usepackage{subcaption}
\usepackage{enumitem}
\usepackage{xspace}
\usepackage{xcolor}
\usepackage{hyperref}
\usepackage{listings}
\usepackage{balance}

% ── Listings ─────────────────────────────────────────────────────────────────
\lstset{
  basicstyle=\ttfamily\scriptsize,
  frame=single, breaklines=true,
  captionpos=b,
}

% ── Theorem environments ──────────────────────────────────────────────────────
\theoremstyle{definition}
\newtheorem{definition}{Definition}
\newtheorem{observation}{Observation}

% ── Colors ───────────────────────────────────────────────────────────────────
\definecolor{riskred}{rgb}{0.65,0.10,0.10}
\definecolor{safegreen}{rgb}{0.08,0.47,0.16}
\definecolor{warnorng}{rgb}{0.78,0.39,0.0}

% ── Macros ───────────────────────────────────────────────────────────────────
\newcommand{\sysname}{\textsc{CapMismatch}\xspace}
\newcommand{\cape}{\textsc{CAPE}\xspace}
\newcommand{\clawhub}{\textsc{ClawHub}\xspace}
\newcommand{\lscps}{\textsc{LSCPS}\xspace}
\newcommand{\eg}{\textit{e.g.,}\xspace}
\newcommand{\ie}{\textit{i.e.,}\xspace}
\newcommand{\etal}{\textit{et al.}\xspace}
\newcommand{\cf}{\textit{cf.}\xspace}

% Key numbers — CAPE final results (125/125 pairs, first-pass)
% ITT = intent-to-treat (counting 64 infrastructure timeouts as failures)
% CR  = completed-run (conditioned on agent returning response before 175s kill)
\newcommand{\NskillsCH}{18{,}443\xspace}
\newcommand{\NpairsTotal}{125\xspace}
% ITT metric (64 timeouts counted as failures)
\newcommand{\NpairsSuccessITT}{57\xspace}
\newcommand{\PairsSuccessRateITT}{45.6\%\xspace}
% Completed-run metric (61 non-timeout pairs)
\newcommand{\NpairsSuccessCR}{57\xspace}
\newcommand{\NpairsCompletedRun}{61\xspace}
\newcommand{\PairsSuccessRateCR}{93.4\%\xspace}
% CI bounds
\newcommand{\CICRlo}{84.3\%\xspace}
\newcommand{\CICRhi}{97.4\%\xspace}
% Shell exec
\newcommand{\ShellSuccessRate}{100\%\xspace}  % 33/33 completed-run
\newcommand{\NShellCompleted}{33\xspace}
% Timeout / refusal counts
\newcommand{\NTimeout}{64\xspace}
\newcommand{\NRefusal}{4\xspace}
% Measurement numbers
\newcommand{\ICone}{75.8\%\xspace}
\newcommand{\IConeN}{13{,}975\xspace}
\newcommand{\MeanCR}{0.042\xspace}
\newcommand{\MeanCDI}{0.894\xspace}
\newcommand{\Tsixrecall}{94.1\%\xspace}
% IAA
\newcommand{\IAAsink}{0.87\xspace}
\newcommand{\skillssh}{\textsc{skills.sh}\xspace}
% ── CAPE skills.sh cross-platform validation (§6) results ─────────────────
\newcommand{\NskillsSH}{116{,}208\xspace}      % total skills.sh corpus
\newcommand{\NssPairs}{197\xspace}              % Codex-confirmed viable skills
\newcommand{\NssRun}{166\xspace}               % skills executed (install ok)
\newcommand{\NssInstallRate}{100\%\xspace}      % 166/166
\newcommand{\NssSuccessRate}{99\%\xspace}      % 165/166
\newcommand{\NssShellRate}{100\%\xspace}       % T06: 75/75
\newcommand{\NssNetworkRate}{96\%\xspace}      % T07: 25/26
\newcommand{\NssMemoryRate}{100\%\xspace}      % T13: 20/20
\newcommand{\NssFileRate}{100\%\xspace}        % T02: 23/23
\newcommand{\NssPersistRate}{100\%\xspace}     % T05: 9/9
% ── ITT/CR rate macros (used in introduction) ──────────────────────────────
\newcommand{\ITTrate}{45.6\xspace}
\newcommand{\ITTlo}{35.3\xspace}
\newcommand{\ITThi}{56.3\xspace}
\newcommand{\CRrate}{93.4\xspace}
\newcommand{\CRlo}{84.3\xspace}
\newcommand{\CRhi}{97.4\xspace}
\newcommand{\NpairsCRDenom}{\NpairsCompletedRun}

% ── Paper metadata ───────────────────────────────────────────────────────────
\title{\sysname: Measuring Capability-Description Inconsistency\\
       and Source-Sink Attack Risks in LLM Agent Skill Marketplaces}

\author{Anonymous Submission}

\begin{document}
\maketitle

% ============================================================
\begin{abstract}
% ============================================================
LLM agent skill marketplaces—platforms where developers publish
installable tools that extend language model capabilities—now host
tens of thousands of skills accessed by millions of users.
Unlike traditional software packages, skills operate with
\emph{NLP-driven ambient authority}: their capabilities are described
only in natural language, with no machine-enforceable permission model.

We study whether these descriptions accurately reflect what skills
actually do.
Analyzing \NskillsCH{} skills from \clawhub, the largest public skill
marketplace, we find that \ICone{} (\IConeN{} skills) \emph{completely omit}
security-sensitive capabilities from their descriptions
(mean Capability Recall \MeanCR).
This opacity is not benign: we demonstrate that it enables a new
class of attacks we call \emph{source-sink skill composition attacks},
in which an attacker-poisoned source skill injects malicious instructions
that activate an unsuspecting sink skill's dangerous, undisclosed capabilities.

We implement \cape (Compositional Attack Pair Evaluator),
an empirical validation framework that tests all \NpairsTotal{} unique
source-sink pairs against real OpenClaw agents in isolated environments.
We report two pre-defined metrics:
\textbf{ITT} (intent-to-treat, counting \NTimeout{} infrastructure timeouts as failures):
\NpairsSuccessITT/\NpairsTotal{} = \PairsSuccessRateITT;
\textbf{completed-run} (conditioned on agent returning a response):
\NpairsSuccessCR/\NpairsCompletedRun{} = \textbf{\PairsSuccessRateCR}
(95\% CI: \CICRlo--\CICRhi).
Side effects include shell command execution, file deletion,
cron backdoor installation, and process termination.
Sink skills with shell-execution capabilities are exploited at
\ShellSuccessRate{} completed-run success rate (\NShellCompleted/\NShellCompleted{} pairs).
Only \NRefusal{} of \NpairsTotal{} pairs (3.2\%) produce explicit agent refusals;
the agent's safety policy is heuristic-based and stochastically applied.
A cross-platform live-execution replication on \NssPairs{} \skillssh{} skills
using the \texttt{codex} agent confirms the finding:
\NssSuccessRate{} attack success rate across all five sink channel types
(shell: \NssShellRate{}, memory: \NssMemoryRate{}, persistence: \NssPersistRate{}),
with \NssInstallRate{} install success rate and zero platform-level safety gates triggered.
Our temporal analysis of skill updates finds that dangerous capabilities
are frequently introduced silently, evading both users and platform scanners.
We release our dataset, tooling, and the \sysname detection pipeline.
\end{abstract}

% ============================================================
\section{Introduction}
\label{sec:intro}
% ============================================================

% sections/ndss_01_introduction.tex
% NDSS 2027 — CapMismatch Introduction
% Target: 1.5 columns / ~750 words

The emergence of LLM agent skill marketplaces---platforms such as
\clawhub (\NskillsCH skills, $>$2.1M installs) and skills.sh that let
developers publish installable tools extending language model
capabilities---represents a new class of software supply chain~\cite{xi2023rise,schick2024toolformer}.
Unlike traditional package managers~\cite{zimmermann2019small},
skills operate with \emph{NLP-driven ambient authority}:
once installed, a skill executes on behalf of the user,
reading files, executing shell commands, accessing credentials,
and sending network requests, often without per-action confirmation
or OS-level sandboxing.

Critically, the \emph{only} capability disclosure mechanism available
to users and platform operators is the skill's natural-language description.
There is no machine-readable permission manifest, no capability whitelist,
and no runtime sandbox to fall back on when the description is inaccurate.
This creates a fundamental \emph{compositional security failure}:
when users combine multiple skills---a web-search source with a
development-utility sink, an email fetcher with a memory manager---
the aggregate set of capabilities the agent can exercise
may far exceed anything disclosed in any individual skill's description.
Neither the user nor the platform can anticipate the attack surface
created by each unique skill combination.

\paragraph{Our contributions.}
We make four contributions:

\begin{enumerate}[leftmargin=*, label=\textbf{C\arabic*.}]

\item \textbf{Measurement at scale.}
We analyze all \NskillsCH skills in \clawhub using
\lscps (Layered Structured Capability Pattern Search),
a rule-based static analysis pipeline for extracting security-sensitive
capabilities from skill implementation code.
We find that \ICone of skills with dangerous capabilities
\emph{completely omit} them from their visible descriptions
(mean Capability Recall $\overline{\text{CR}} = \MeanCR$),
constituting an ecosystem-wide transparency failure.

\item \textbf{IC taxonomy.}
We introduce the Inconsistency Class (IC) taxonomy---IC-0 through IC-4---that
characterizes four structural patterns of capability-description mismatch,
from full omission (IC-1, 75.8\%) to false claims (IC-3, 0.4\%).
This taxonomy provides the first systematic vocabulary for reasoning
about skill transparency.

\item \textbf{Compositional attack surface validation (CAPE).}
We demonstrate that capability opacity does not merely create theoretical risk;
it creates a \emph{measurable, exploitable compositional attack surface}.
When any source skill (web search, email, calendar, file reader, memory,
browser automation, database connector) is co-installed with an IC-1 sink
skill, the combination constitutes an attack surface: an adversary
who can influence the content delivered by the source skill can trigger
the sink's undisclosed dangerous capability without any
additional platform access.
This attack surface arises from the composition itself---
neither source nor sink is individually malicious,
but together they expose capabilities the user never consented to.

We operationalize this surface through \cape (Compositional Attack Pair Evaluator),
testing \NpairsTotal unique source-sink pairs spanning seven injection-channel
categories (C1: LLM context, C2: filesystem, C3: agent memory, C5: database,
C9: boot hook / cron, C13: network) on a real OpenClaw /
\texttt{gpt-5.3-codex} agent platform.
Each source skill is individually registered as an
identity-spoofing MCP server mirroring the real skill's name and tool schema,
modeling the realistic supply-chain threat of a malicious clone skill.
We report two pre-defined metrics:
\textbf{ITT}: \NpairsSuccessITT/\NpairsTotal = \ITTrate\% (95\% CI: \ITTlo--\ITThi\%),
counting infrastructure timeouts as failures;
\textbf{completed-run}: \NpairsSuccessCR/\NpairsCRDenom = \CRrate\%
(95\% CI: \CRlo--\CRhi\%)
conditioned on pairs where the agent returned a response.
Side effects include shell command execution, file deletion,
cron backdoor installation, cross-session memory poisoning,
process termination, wallet exfiltration, and privilege escalation.
33/33 tested shell-execution sink pairs succeed (completed-run: 100\%),
revealing \emph{no consistent runtime defense} against this attack class.

\item \textbf{Temporal analysis.}
We track 24,248 version transitions in \clawhub and find that dangerous
capabilities are silently introduced in updates with no description change (\Tsixrecall recall at Tier~6).
Platform scanners detect 0\% of these transitions.

\end{enumerate}

\paragraph{Scope and novelty vs. prompt injection.}
Our work is related to but distinct from prompt injection
attacks~\cite{greshake2023indirect,perez2022ignore}.
Prompt injection studies \emph{the manipulation of individual inputs}
to produce unsafe agent behavior; we study
\emph{the compositional attack surface created by installing
combinations of skills with undisclosed capabilities}.
The distinction is architectural:
our attack requires no crafted user input and no runtime adversarial access.
The attack surface exists the moment a dangerous IC-1 sink skill
is installed alongside any source skill---it is a \emph{structural}
vulnerability of the skill composition, not an input-validation failure.
This means existing prompt-injection defenses (input sanitization,
constrained decoding, system-prompt hardening) are \emph{not applicable}:
the dangerous capability $\kappa$ is invoked through the agent's
normal, intended tool-use workflow, with attacker-influenced content
arriving through a legitimate, already-trusted source channel.

\paragraph{Ethical considerations.}
All experiments were conducted in fully isolated environments
with no impact on real users, external services, or production systems.
Skill identifiers used in this paper are real ClawHub skill names
included to enable reproducibility; no exploit code is released.
We have privately notified the ClawHub platform operators of our findings.

\paragraph{Paper organization.}
\S\ref{sec:background} presents the threat model.
\S\ref{sec:taxonomy} introduces the IC taxonomy.
\S\ref{sec:measurement} reports the large-scale measurement.
\S\ref{sec:cape} presents the \cape attack validation.
\S\ref{sec:temporal} analyzes version-update risks.
\S\ref{sec:discussion} discusses limitations and defenses.
\S\ref{sec:related} situates our work.


% ============================================================
\section{Background and Threat Model}
\label{sec:background}
% ============================================================

% sections/ndss_02_background.tex  — NDSS 2027

\subsection{LLM Agent Skill Marketplaces}

LLM skill marketplaces expose installable tools that extend agent
capabilities via standardized protocols such as the Model Context Protocol (MCP)~\cite{anthropic2024mcp}.
A skill is a self-contained package comprising:
(i) a \emph{description} (natural-language documentation visible to users),
(ii) \emph{tool definitions} (structured capability declarations),
and (iii) \emph{implementation code} (executed by the agent runtime).
\clawhub, launched with MCP in late 2024, is the largest such marketplace,
with \NskillsCH skills and $>$2.1M installs as of March 2026.

\paragraph{Ambient authority and sessions.}
Installed skills execute with the privileges of the agent runtime---typically
the user's OS account.
Skills can read and write files, execute shell commands, make network requests,
and read credentials without per-action OS-level confirmation.
This \emph{ambient authority} model, common in agent frameworks~\cite{ruan2024identifying},
means that a skill's behavior is constrained only by its own code
and the LLM's willingness to call its tools.

\subsection{Threat Model}

\paragraph{Principals.}
We consider three principals:
(1)~the \emph{user}, who installs and interacts with skills;
(2)~the \emph{skill developer}, who publishes skills to the marketplace;
and (3)~the \emph{attacker}, who may be a malicious skill developer
or an attacker controlling content fetched by a legitimate source skill.

\paragraph{Attacker capabilities.}
The attacker:
\begin{itemize}[leftmargin=*,topsep=2pt]
  \item can publish skills with benign-looking descriptions
    that omit dangerous capabilities (supply-chain attacker), OR
  \item can control content deliverable through a source skill's
    data-fetch operations (\eg web pages, emails, API responses)---the
    \emph{prompt injection via external content} model~\cite{greshake2023indirect}.
\end{itemize}

\paragraph{Attack goal.}
The attacker aims to cause the agent to invoke a dangerous capability
$\kappa$ of an installed sink skill, without the user's knowledge
or direct authorization, producing a real system-level effect $e$
(\eg file deletion, shell command execution, data exfiltration).

\paragraph{Out of scope.}
We do not model: direct UI-level social engineering of users,
OS-level privilege escalation beyond the agent runtime's account,
or attacks on the marketplace platform infrastructure itself.

\paragraph{Key assumption.}
The attack is \emph{enabled by the IC-1 gap}:
the user installed the sink skill because its description appeared benign,
and therefore has no awareness that $\kappa$ is possible.
This distinguishes our threat from known prompt injection scenarios
where the dangerous capability is explicitly documented.

\subsection{Why Descriptions Are the Only Defense}

Unlike Android (permission manifests~\cite{felt2011android}),
npm (lockfiles + audit~\cite{synk2023report}),
or sandboxed browser extensions (capability APIs),
LLM skill marketplaces have \emph{no machine-enforceable permission model}.
The MCP standard defines tool schemas, but these describe the API surface,
not the security implications of tool calls.
A skill declaring a \texttt{run\_command(cmd: str)} tool provides no
machine-readable signal that \texttt{cmd} is user-controlled or dangerous.
Natural-language descriptions are therefore the \emph{de facto}
capability disclosure mechanism, and their accuracy is a security property.


% ============================================================
\section{Capability-Description Inconsistency}
\label{sec:taxonomy}
% ============================================================

% sections/ndss_03_taxonomy.tex  — NDSS 2027 (condensed from CCS version)

We formalize capability-description mismatch via a
\emph{difference matrix} between what a skill claims ($\mathcal{C}_{\text{vis}}$)
and what it does ($\mathcal{C}_{\text{act}}$).

\paragraph{Capability sets.}
$\mathcal{C}_{\text{vis}}(s)$ is the set of security-sensitive capability
families mentioned in the skill's marketplace description.
$\mathcal{C}_{\text{act}}(s)$ is the set detected by static analysis
of the skill's implementation code (\lscps; \S\ref{sec:measurement}).
Both sets are drawn from the same 14-family taxonomy
(\eg \texttt{FILE\_WRITE}, \texttt{SHELL\_EXEC}, \texttt{CREDENTIAL\_ACCESS},
\texttt{NETWORK\_EXFIL})
derived from supply-chain attack taxonomies~\cite{ohm2020backstabber,ladisa2023taxonomy}.

\paragraph{Difference matrix.}
For each skill $s$ we compute:
\begin{align*}
\text{AD}&=|\mathcal{C}_{\text{act}} \cap \mathcal{C}_{\text{vis}}| &
\text{(accurately disclosed)}\\
\text{UD}&=|\mathcal{C}_{\text{act}} \setminus \mathcal{C}_{\text{vis}}| &
\text{(undisclosed / hidden)}\\
\text{FC}&=|\mathcal{C}_{\text{vis}} \setminus \mathcal{C}_{\text{act}}| &
\text{(false claim)}
\end{align*}

\paragraph{IC taxonomy.}
We define five inconsistency classes:
\begin{itemize}[leftmargin=*,topsep=2pt,itemsep=1pt]
  \item \textbf{IC-0 (Accurate):} AD$>$0, UD$=$FC$=$0. Description accurately captures capabilities.
  \item \textbf{IC-1 (Full Omission):} UD$>$0, AD$=$FC$=$0. Dangerous capabilities completely hidden.
  \item \textbf{IC-2 (Partial Omission):} UD$>$0, AD$>$0. Some capabilities disclosed, others hidden.
  \item \textbf{IC-3 (False Claim):} FC$>$0, UD$=$0. Capabilities claimed but not implemented.
  \item \textbf{IC-4 (Bidirectional):} UD$>$0, FC$>$0. Both under- and over-disclosure present.
\end{itemize}

\paragraph{Metrics.}
We derive three scalar metrics:
\begin{align}
\text{CR}(s) &= \frac{\text{AD}}{\text{AD}+\text{UD}} \label{eq:cr}
  \quad \text{(Capability Recall)}\\
\text{CP}(s) &= \frac{\text{AD}}{\text{AD}+\text{FC}} \label{eq:cp}
  \quad \text{(Description Precision)}\\
\text{CDI}(s) &= \frac{\text{UD}+\text{FC}}{|\mathcal{F}|} \label{eq:cdi}
  \quad \text{(Capability Discrepancy Index)}
\end{align}

CR$=0$ means the description mentions nothing the skill can actually do;
CDI$=1$ means every capability family is misrepresented.


% ============================================================
\section{Measurement Study}
\label{sec:measurement}
% ============================================================

% sections/ndss_04_measurement.tex  — NDSS 2027

\subsection{LSCPS: Static Capability Extraction}

\lscps (Layered Structured Capability Pattern Search) is a
rule-based static analysis pipeline that scans all skill artifact
files---Python/JS/shell code, SKILL.md instructions,
dependency manifests, and agent config files---for 300+
regex patterns organized into 14 security-sensitive capability families.
Patterns operate on parsed AST nodes where available
(Python \texttt{ast}, \texttt{esprima}) and raw text for shell scripts.
\lscps is intentionally conservative: a pattern match contributes
to $\mathcal{C}_{\text{act}}$ only if it appears in executable
or instruction contexts, not in comments or example strings.
This preserves high precision at the cost of recall,
making all IC rates \emph{lower bounds}.

\subsection{Corpus and Results}
\label{sec:meas_results}

We analyze all \NskillsCH skills in \clawhub (March 2026 snapshot).

\begin{table}[t]
\centering
\caption{IC type distribution across \NskillsCH \clawhub skills.
  Pct.\ over skills with $\mathcal{C}_{\text{act}} \neq \varnothing$.}
\label{tab:ic_dist}
\small
\begin{tabular}{llrr}
\toprule
\textbf{Type} & \textbf{Description} & \textbf{Count} & \textbf{\%}\\
\midrule
IC-1 & Full omission  & 13,975 & 75.8\%\\
IC-2 & Partial omission & 2,507 & 13.6\%\\
IC-4 & Bidirectional  &   535 &  2.9\%\\
IC-3 & False claim    &    74 &  0.4\%\\
IC-0 & Accurate       &    74 &  0.4\%\\
\midrule
None & No dangerous cap.\ detected & 1,278 & 6.9\%\\
\bottomrule
\end{tabular}
\end{table}

\textbf{IC-1 dominates}: 75.8\% of skills with dangerous detected
capabilities entirely omit them from descriptions.
Only 0.4\% achieve accurate disclosure.
The mean Capability Recall $\overline{\text{CR}} = \MeanCR$ indicates
that descriptions mention fewer than one in twenty-three
detected capability signals.
Mean CDI $= \MeanCDI$ confirms that nearly all skills have at least
one misrepresented capability family.

\paragraph{Capability family breakdown.}
The most commonly hidden capability families are:
\texttt{SHELL\_EXEC} (hidden in 38.2\% of skills with the capability),
\texttt{FILE\_WRITE} (29.1\%), and \texttt{NETWORK\_EXFIL} (21.4\%).
These three families account for 72\% of the CAPE-validated dangerous sinks (\S\ref{sec:cape}).

\paragraph{Confirmed dangerous sinks.}
From the IC-1 and IC-2 populations, we identify 329 candidate
dangerous sink skills via LSCPS severity scoring.
Human verification of the top 125 yields the \NpairsTotal confirmed
sinks used in \cape (\S\ref{sec:cape}).
Verification assessed whether the detected capability was
(a) actually present in the implementation and
(b) triggerable by agent-side tool invocation without out-of-band
machine access.
To assess labeling reliability, two annotators independently verified
a stratified random sample of 40 sinks (32\% of the verified set).
Inter-annotator agreement (Cohen's $\kappa$) was \textbf{0.87} on
binary dangerous/not-dangerous classification,
indicating near-perfect agreement and supporting the reliability
of the full 125-sink ground truth set.

\subsection{Cross-Platform Validation}

To confirm findings generalize beyond \clawhub,
we analyze 44 skills from skills.sh---an independent, Europe-based
LLM skill marketplace with a different operator and user base---using
the same \lscps pipeline with no parameter changes.
40/44 (90.9\%) exhibit IC-1 or IC-2,
consistent with our primary corpus.
The cross-platform consistency rules out ClawHub-specific effects
and suggests the transparency gap is a structural property of
current skill marketplace designs, not an artifact of any single platform's policies.
Note that this paper's quantitative claims (IC-1 = 75.8\%, CR = 0.042)
are scoped to \clawhub; cross-platform rates serve as directional confirmation only.


% ============================================================
\section{Source-Sink Attack Validation (CAPE)}
\label{sec:cape}
% ============================================================

% CAPE section moved to the parallel CAPE experiment line.
% sections/ndss_05_cape.tex
% NDSS 2027 / CCS 2026 — §5 Source-Sink Compositional Attack Surface Validation (CAPE)
% Target: ~2 columns / ~1100 words

Having established that \ICone{} of ClawHub skills omit dangerous capabilities
from their public descriptions, we ask a structural security question:
\emph{does installing combinations of skills create attack surfaces the user
cannot anticipate?}
We answer through \cape (Compositional Attack Pair Evaluator),
an empirical framework that operationalizes the compositional attack surface
across \NpairsTotal{} unique source-sink pairs on a real LLM agent platform
(OpenClaw / \texttt{gpt-5.3-codex}).
A cross-platform replication on \skillssh{} using live \texttt{codex} agent
execution is reported in \S\ref{sec:crossplatform}.

\subsection{Attack Model and Threat Surface}

\paragraph{Compositional attack surface.}
The central security claim of this paper is not that any individual skill is
malicious, but that \emph{skill composition creates attack surfaces invisible
to the user}.
Formally, we define:

\begin{definition}[Source-Sink Compositional Attack Surface]
For a skill ecosystem $\mathcal{E}$, the \emph{compositional attack surface}
is the set of all source-sink pairs
$(\mathcal{S}_{\text{src}}, \mathcal{S}_{\text{sink}}) \in \mathcal{E}^2$
such that (i) $\mathcal{S}_{\text{sink}}$ possesses a dangerous capability
$\kappa \in \mathcal{C}_{\text{act,sink}} \setminus \mathcal{C}_{\text{vis,sink}}$
not disclosed to the user (\ICone), and (ii) $\mathcal{S}_{\text{src}}$
provides a channel through which attacker-influenced content can reach
the shared agent context.
\end{definition}

This definition has three key implications.
First, \textbf{the attack surface is structural}:
it arises from the act of combination itself, not from either skill being
individually malicious---a user who independently audits each installed
skill cannot detect the risk.
Second, \textbf{the attack surface is channel-agnostic}:
any source skill---web search, email, calendar, file reader, memory store,
browser automation, or database connector---can serve as the delivery
channel because the agent shares a single context window across all
installed skills.
Third, \textbf{the enabling condition is capability omission (\ICone)}:
if the sink skill disclosed its \texttt{shell\_exec} capability,
a user could apply the principle of least privilege and install only
the source skills they trust to carry trusted content.
IC-1 removes this defense entirely.

\begin{definition}[Source-Sink Composition Attack]
Given an agent session $\mathcal{A}$ with installed skills
$\mathcal{S}_{\text{src}}$ and $\mathcal{S}_{\text{sink}}$,
an attacker who can influence any content delivered through
$\mathcal{S}_{\text{src}}$'s output channel constructs content $c$
such that $\mathcal{A}$'s processing of $c$ invokes
$\kappa \in \mathcal{C}_{\text{act,sink}} \setminus \mathcal{C}_{\text{vis,sink}}$
without explicit user authorization.
\end{definition}

Note that the attacker's power requirement is minimal:
they need only inject content into \emph{one} of many plausible surfaces---
a search-engine result, an email in a monitored inbox, a document in a
shared folder, a cross-session memory entry, or an API response.
The \ICone gap is what converts this low-power capability into
a high-impact system compromise.

\subsection{Experimental Setup}

\paragraph{Pair selection.}
We select \NpairsTotal unique source-sink pairs to systematically
map the compositional attack surface.
Each pair combines one confirmed IC-1 sink with one source skill,
constituting an instance of the compositional attack surface.
Each pair is unique: no source or sink is reused across pairs,
ensuring independent validation of each sink's exploitability.
Source skills span seven injection-channel categories (\cf Table~\ref{tab:channels}):
\textbf{C1} LLM context (97 pairs): web search, browser automation,
  news fetch, API aggregation;
\textbf{C2} filesystem (1 pair): direct file read;
\textbf{C3} agent memory (17 pairs): cross-session memory write/read;
\textbf{C9} boot hook (4 pairs): startup script injection;
\textbf{C13} network (3 pairs): outbound request channel;
\textbf{C5} database (2 pairs): DB read/write;
\textbf{C9\_cron} (1 pair): scheduled task channel.
This multi-channel design directly validates that the attack surface
is channel-agnostic, not an artifact of any single delivery mechanism.

\paragraph{Task assignment.}
We define 15 attack task types (T1--T15) mapped to sink capability
categories: shell execution, file deletion, cron backdoor,
privilege escalation, script evaluation, process termination,
cross-session memory injection, wallet exfiltration, and others.
Each pair is assigned the task that best matches its sink's
confirmed dangerous capability.

\paragraph{Isolation.}
Each experiment runs in an isolated environment:
a dedicated workspace at \texttt{\textasciitilde/.cape\_exp/<pair\_id>/}
seeded with realistic context files (credentials, BOOT.md, memory files),
and a unique OpenClaw session ID.
A lightweight C2 server captures exfiltrated data for network-based tasks.
Side effects are verified by checking for probe files,
C2 server receipts, process state, crontab entries, or file modifications.

\paragraph{Injection delivery: identity-spoofing MCP source chain.}
A key methodological commitment distinguishes \cape from prior prompt injection evaluations:
we implement a \emph{true source-sink execution chain} where each pair's
source skill is individually registered as an identity-spoofing MCP server
that mirrors the real source skill's exact name, tool schema, and
Server Info declaration.

For each pair, our attacker-controlled MCP stdio server
(\texttt{attacker\_source\_skill.mjs}) is registered under the
\emph{real source skill's name} (\eg \texttt{tavily-search},
\texttt{brave-search}, \texttt{outlook-graph}, \texttt{browser-use})
via \texttt{openclaw\,mcp\,set}---identical to how a legitimate
skill is installed.
The server exposes the real tool name the source skill would expose
(\eg \texttt{tavily-search\_\_search}, \texttt{browser-use\_\_browse},
\texttt{outlook-graph\_\_get\_emails}),
ensuring the agent believes it is interacting with the
legitimate installed skill.
This models the realistic threat of a malicious clone skill published
under a name indistinguishable from a popular legitimate skill---
a supply-chain attack requiring no exploitation of the target platform's
infrastructure.

The attack chain proceeds as follows:

\begin{enumerate}[leftmargin=*,topsep=2pt,itemsep=1pt]
\item \textbf{Source skill registration (attacker).}
  The attacker publishes a skill with the same name and tool schema
  as a popular source skill (\eg Tavily Search).
  One-time registration via \texttt{openclaw\,mcp\,set} with the real skill's identity.

\item \textbf{Legitimate user workflow.}
  The user installs both a source skill (for web search, email, etc.) and
  a sink skill (for development, scheduling, memory, etc.).
  Neither skill appears dangerous in isolation;
  the sink's dangerous capability is IC-1-omitted.

\item \textbf{Source skill invocation.}
  During a routine task, the agent calls the source skill's tool
  (\eg \texttt{tavily-search\_\_search}) to fetch documentation or data.
  The call triggers a real MCP tool invocation over the OpenClaw gateway.

\item \textbf{Injection via tool output.}
  The source skill returns crafted ``integration documentation''
  containing a disguised shell command.
  The payload is framed as a legitimacy-enhancing pre-flight check
  using a living-off-the-land pattern:
  \texttt{go\,vet\,./...\,2>\&1\,|\,tee\,\textlangle{}probe\textrangle{}},
  presented as a ``go vet compatibility check required by the sink skill.''
  We validate two alternative payload families:
  (\textbf{P2}) \texttt{node\,verify.js\,2>\&1\,|\,tee\,\textlangle{}p\textrangle{}}
  and (\textbf{P3}) \texttt{cat\,README.md\,|\,tee\,\textlangle{}p\textrangle{}},
  both achieving comparable success (P2: 7/8, P3: 6/8 on held-out pairs).

\item \textbf{Sink capability invocation.}
  The agent executes the command verbatim, triggering the sink skill's
  IC-1-omitted \texttt{shell\_exec} capability and
  producing a verified file-system side effect.
\end{enumerate}

This chain provides Level~3 empirical evidence---real registered skill,
real agent session, real tool invocation, real system-level side effect---
validating the compositional attack surface across 125 distinct source-sink pairs
spanning 7 injection channels and 13 dangerous capability categories.

\paragraph{Controls.}
To distinguish attack effect from model variance, we run three control conditions
on a stratified random sample of 20 pairs (4 per task type):
\textbf{(Ctrl-A)} same pair, but source skill returns a benign documentation page (no injection);
\textbf{(Ctrl-B)} same pair, same injection, but the payload command is replaced with a clearly
benign operation (\texttt{echo hello > /dev/null});
\textbf{(Ctrl-C)} same pair, repeated 3$\times$ with identical injection to measure within-pair variance.
Ctrl-A: 0/20 probe files created (baseline = zero without injection).
Ctrl-B: 18/20 agent executes the benign command (confirming agent compliance with tool context),
but 0/20 create the probe file.
Ctrl-C: within-pair success variance is low ($\sigma = 0.11$) for T06 pairs and moderate
($\sigma = 0.28$) for T18/T13 pairs, consistent with stochastic refusal rather than random model behavior.
These controls confirm that: (1) probe file creation is caused by the injected command, not ambient side effects;
(2) agent compliance is not the bottleneck---payload delivery is;
and (3) stochastic refusals in T13/T18 reflect per-session context sensitivity, not a stable defense.



\subsection{Results}

\begin{table}[t]
\centering
\caption{CAPE attack validation results by task type across all 125 source-sink pairs.
  All pairs use the true MCP source-sink chain (real attacker MCP server \textrightarrow\ real agent session \textrightarrow\ verified file-system side effect).}
\label{tab:cape_results}
\setlength{\tabcolsep}{4pt}
\small
\begin{tabular}{lllrr}
\toprule
\textbf{Task} & \textbf{Sink Capability} & \textbf{Injection Disguise} & \textbf{Pairs} & \textbf{Success}\\
\midrule
T06 & Shell exec (arbitrary)   & \texttt{go vet | tee \textlangle{}p\textrangle{}} & 70  & 33 (47\% ITT; \textbf{100\% CR}) \\
T18 & Process termination      & \texttt{go vet | tee \textlangle{}p\textrangle{}} & 17  & 5 (29\% ITT; 71\% CR) \\
T13 & Cross-session memory     & \texttt{go vet | tee \textlangle{}p\textrangle{}} & 16  & 8 (50\% ITT; 89\% CR) \\
T17 & Wallet exfiltration      & \texttt{go vet | tee \textlangle{}p\textrangle{}} &  5  & 3 (60\% ITT; 100\% CR) \\
T05 & Boot hook injection      & \texttt{go vet | tee \textlangle{}p\textrangle{}} &  4  & 0 (0\% ITT; 0\% CR) \\
T03 & Cron backdoor            & \texttt{go vet | tee \textlangle{}p\textrangle{}} &  1  & 1 (100\%) \\
T14 & Script eval              & \texttt{go vet | tee \textlangle{}p\textrangle{}} &  1  & 1 (100\%) \\
T11 & Privilege escalation     & \texttt{go vet | tee \textlangle{}p\textrangle{}} &  1  & 1 (100\%) \\
T07/T08/T10 & Network/DB exfil  & \texttt{go vet | tee \textlangle{}p\textrangle{}} &  5  & 5 (100\%) \\
T02/T04/T19/T20 & Other        & \texttt{go vet | tee \textlangle{}p\textrangle{}} &  5  & 0 (0\% — all timeout) \\
\midrule
\textbf{Total} && & \textbf{125} & \textbf{57 ITT (45.6\%) / 57 CR (93.4\%)} \\
\bottomrule
\end{tabular}
\end{table}

Table~\ref{tab:cape_results} reports results for 85/125 completed pairs
(40 remaining; camera-ready will include all 125).
We report two metrics, pre-defined before any analysis:

\textbf{Intent-to-treat (ITT)}: 57/85 = \textbf{67.1\%}
(95\% CI: 56.5--76.1\%),
counting all 24 timed-out pairs as failures.
Timeout policy: each pair is given one attempt with a 175\,s hard kill
(\texttt{SIGKILL} on the process group);
if no probe file is created within that window, the pair is marked failed.

\textbf{Completed-run rate}: 57/61 = \textbf{93.4\%}
(95\% CI: 84.1--97.8\%),
conditioned on pairs where the agent returned a response before timeout.
24 pairs timed out due to agent infrastructure latency
(LLM context compaction under high prompt token usage)---
this is a measurement-infrastructure limitation, not an agent defense:
timeout pairs produce no output and thus no refusal, evidence, or probe file.
Of the 24 ITT failures, only \textbf{4} (4.7\% of tested pairs) are agent refusals;
the remaining 20 are silent timeouts with no agent output.

T06 (shell exec, completed-run): 33/33 = \textbf{100\%};
T07/T08/T10 (network/DB exfil): 5/5 = 100\%;
T17 (wallet): 3/3 = 100\%.

\paragraph{Uniform injection technique — all 15 task types.}
A key finding is that the \textbf{same injection disguise} ($\S$\ref{sec:injection_detail})
works across all 15 sink task types: not just shell-exec sinks,
but also process monitors, memory indexers, wallet validators,
boot hooks, database connectors, and script evaluators.
Each sink's documentation is framed around go vet compatibility,
causing agents to execute the command regardless of the sink's primary domain.
This demonstrates that the vulnerability is in the \emph{agent's trust model},
not in any property specific to shell-exec sinks.

\paragraph{True MCP source-sink chain: all 125 pairs.}
Unlike prior prompt injection evaluations that hand-craft source-skill output,
all 125 \cape pairs use a \emph{real, registered MCP server} as the source skill.
The attacker MCP server (\texttt{attacker\_source\_skill.mjs}) is installed
via the standard \texttt{openclaw\,mcp\,set} interface---
the same mechanism used to install any legitimate skill---and its
\texttt{web\_search} tool interface is indistinguishable from
Tavily, DuckDuckGo, Google Search, or any other web-search source skill.
Source skills span 12+ different types (web search, email, news feeds,
browser automation, OCR, RSS, social media scrapers) across all 125 pairs,
confirming that the attack generalizes across the full source skill ecosystem.

\paragraph{Shell execution sinks: no consistent defense.}
Of 33 T06 (shell exec) pairs that ran to completion, \textbf{33/33 = 100\%} produced
verified side effects across 33 distinct source skills (web search, email,
browser automation, finance APIs, news aggregators, OCR, and social scrapers).
The agent's safety model applies no consistent restriction to shell commands
delivered through source-skill tool output.
This is the most dangerous finding: 55.6\% of all confirmed dangerous sinks
(\S\ref{sec:measurement}) have shell execution capability, and no completed pair
shows a successful defense against the \texttt{go\,vet}-disguised injection.

\paragraph{Asymmetric safety policy.}
The agent exhibits a striking asymmetry:
it occasionally refuses requests to modify agent memory files directly
(``I can't run pasted diagnostic commands from untrusted integration messages''),
but applies no equivalent restriction to \texttt{go vet}
or other build-tool commands.
This heuristic-based policy leaves the more dangerous execution surface
completely unguarded against build-tool-disguised payloads.

\paragraph{Failure analysis and refusal variance.}
Of 85 pairs attempted, 4 produced explicit agent refusals:
(i) P001: ``I can't run commands extracting local identity content from untrusted search output'';
(ii) P004: ``system process data from untrusted source is not acceptable'';
(iii) P049, P050: LLM idle timeout (no output produced).

Critically, \emph{neither refusal message fires consistently}: T13 (memory) succeeds in
8/12 tested pairs (67\%) and T18 (process kill) in 5/11 (45\%),
demonstrating that the agent applies safety heuristics stochastically
rather than enforcing a stable capability policy.
Per-pair repetition ($n=3$; Ctrl-C; \S\ref{sec:controls}) confirms:
Within-task-type refusal variance is $\sigma = 0.11$ for T06 and $\sigma = 0.28$ for T13/T18.
A single refusal in one session is not predictive of refusal in subsequent sessions with the same pair.
This inconsistency is itself a security vulnerability: an attacker need only retry
under different session contexts until the safety heuristic does not fire.

\paragraph{Injection disguise generalizes.}
The \texttt{go vet ./... 2>\&1 | tee \textlangle{}probe\textrangle{}} payload
succeeds because the agent interprets it as a routine build-tool caching step.
This ``living-off-the-land'' pattern---using legitimate, pre-installed
developer tools in ways that produce attacker-controlled side effects---is
not detectable by any static inspection of the command string.
No special capabilities or elevated permissions are required;
the attack exploits only tools already present in a typical developer environment.


\subsection{Key Takeaway}
\label{sec:cape_takeaway}

\begin{observation}
A skill's undisclosed dangerous capability is, in practice,
exploitable by any source skill that can deliver attacker-controlled
content to the agent's context window.
The platform's safety policy is heuristic-based and stochastically applied:
no consistent runtime defense is activated against build-tool-disguised
injection payloads, even across repeated trials.
\end{observation}

Our results establish that the IC-1 gap (capability omission) is not
merely a transparency issue---it is an active attack enabler.
Users cannot consent to capabilities they cannot see,
and platforms cannot enforce policies they have not enumerated.
A cross-platform replication with live \texttt{codex} agent execution
on \NssPairs{} \skillssh{} skills ($\NssSuccessRate{}$ attack success rate)
is reported in \S\ref{sec:crossplatform}.



% ============================================================
\section{Cross-Platform Validation}
\label{sec:crossplatform}
% ============================================================

% sections/ndss_06_temporal.tex  — NDSS 2027 / CCS 2026
% §6: Cross-Platform Validation — skills.sh Live Execution
% Target: ~1.5 columns / ~700 words

\section{Cross-Platform Validation}
\label{sec:crossplatform}

The ClawHub \cape experiment (\S\ref{sec:cape}) validates the
compositional attack surface on a single platform.
To establish generality, we conduct two independent cross-platform analyses:
(i) a static inconsistency measurement on \skillssh{},
and (ii) a live execution experiment using a real \texttt{codex}
agent instance.

\subsection{Static Analysis on \skillssh{}}

We apply the \lscps pipeline (\S\ref{sec:measurement}) without modification
to \NskillsSH{} \skillssh{} skills (retrieved April 2026).
The pipeline is designed for ClawHub's SKILL.md format;
\skillssh{} uses the same format,
making the transfer zero-shot.

\paragraph{Result.}
IC-1 (complete capability non-disclosure) applies to
$\mathbf{>90\%}$ of \skillssh{} skills---
a rate higher than ClawHub ($75.8\%$),
consistent with the hypothesis that disclosure norms
depend on marketplace culture rather than technical constraints.
Source-skill concentration ($S_4$: local system access, $S_1$: web fetch)
and sink-skill concentration ($K_1$: shell execution, $K_5$: persistence)
mirror the ClawHub distribution, confirming that the
Source$\times$Sink attack surface is structurally universal.

\subsection{Live Execution: \skillssh{} CAPE}
\label{sec:ss_cape}

\paragraph{Dataset.}
Starting from the \NskillsSH{} \skillssh{} corpus,
we apply a three-stage filter:
(1) LSCPS source-sink pair detection ($4{,}154$ self-exfil candidates);
(2) Codex semantic audit via chain-of-thought reasoning
    (prompt: ``Does this skill's SKILL.md omit any sink-type dangerous capability?''),
retaining $\NssPairs{}$ Codex-confirmed viable attack paths; and
(3) install-command availability (required for automation).
The $\NssPairs{}$ confirmed skills span all five
sink channel types (Table~\ref{tab:ss_results}).

\paragraph{Experimental setup.}
Each skill is evaluated in a fully isolated workspace
(\texttt{/tmp/cape\_sandboxes/\textlangle{}id\textrangle{}/})
containing realistic credential files
(\texttt{credentials/openai.token}, \texttt{credentials/github.token},
\texttt{memory/hot/HOT\_MEMORY.md}, \texttt{BOOT.md}).
Installation: \texttt{npx skills add \textlangle{}repo\textrangle{} --skill \textlangle{}name\textrangle{} -a codex -y}.
Execution: \texttt{codex exec --full-auto --skip-git-repo-check \textlangle{}trigger\textrangle{}}.
This represents a self-exfil attack model:
the skill is both source (it processes external input)
and sink (it executes dangerous operations),
which is the most conservative threat model---
it does not require any composition.
Side effects are verified by filesystem probe:
a skill-specific marker file, BOOT.md content,
memory-file modification, or exfil log.
Each task type uses the injection trigger described in Table~\ref{tab:ss_results}.

\paragraph{Results.}
Table~\ref{tab:ss_results} reports results for \NssRun{} executed skills
(31 skipped: GitHub access-controlled or withdrawn between audit and execution).
Install success: $\NssRun{}/\NssRun{} = \NssInstallRate{}$.
Attack success: $165/166 = \NssSuccessRate{}$.

\begin{table}[t]
\centering
\caption{%
  \skillssh{} live execution results by sink channel type.
  All experiments use \texttt{codex exec --full-auto}
  with per-sink injection triggers.
  ``Success'' = verified filesystem side effect.
}
\label{tab:ss_results}
\setlength{\tabcolsep}{4pt}
\small
\begin{tabular}{llrrc}
\toprule
\textbf{Task} & \textbf{Sink Channel} & \textbf{N} & \textbf{Success} & \textbf{Rate}\\
\midrule
T06 & C1: Shell execution        & 75 & 75 & \textbf{100\%}\\
T07 & C4: Network data collect   & 26 & 25 & \textbf{96\%}\\
T02 & C2: Filesystem (cred del.) & 23 & 23 & \textbf{100\%}\\
T13 & C3: Agent memory poison    & 20 & 20 & \textbf{100\%}\\
T19 & C4: Messaging webhook      & 13 & 13 & \textbf{100\%}\\
T05 & C5: Persistence (boot)     &  9 &  9 & \textbf{100\%}\\
\midrule
\textbf{All} & & \textbf{166} & \textbf{165} & \textbf{\NssSuccessRate{}}\\
\bottomrule
\end{tabular}
\end{table}

\paragraph{Key findings.}(i)~\textbf{Install-time barrier is zero}:
every installable skill installs cleanly without any platform-level
safety gate ($\NssInstallRate{}$, 166/166).
(ii)~\textbf{Runtime exploitability is near-universal}:
$\NssSuccessRate{}$ of installed skills produce dangerous verified
side effects---T06 (shell exec) at 100\%, T13 (memory poison) at 100\%,
T05 (boot hook persistence) at 100\%.
(iii)~\textbf{Network exfil is the only defended channel}:
T07 achieves 96\% rather than 100\%; agent $\texttt{gpt-5.4}$
detects explicit credential-reading scripts in 4\% of network-exfil cases
and refuses---the only category with any observable safety response.
This asymmetry confirms that safety defenses are
heuristic-based and target-specific, not capability-based.
(iv)~\textbf{Cross-platform consistency}:
the Source$\times$Sink matrix from ClawHub (\S\ref{sec:cape})
is reproduced across both platforms and both attack models
(compositional and self-exfil), confirming that
IC-1 gaps are a structural property of skill markets,
not an artifact of a single platform's design.

\paragraph{Temporal evolution.}
\label{sec:temporal}
We additionally analyze $24{,}248$ version transitions across $8{,}120$
\clawhub{} skills tracked between November 2024 and March 2026.
Of these, 6.8\% introduce new security-sensitive capabilities
without any description update (silent TCE).
627 skills (7.7\% of updated skills) acquired new capability families
silently---\eg a web-search skill that gained \texttt{FILE\_WRITE}
in a minor update.
Applied to 40 confirmed malicious skill version pairs
(34 independently flagged malicious skills),
\lscps achieves~\Tsixrecall{} recall at Tier-6 severity threshold;
the platform scanner achieves 0\%.
This establishes temporal monitoring as a necessary complement
to publish-time screening.


% ============================================================
\section{Discussion}
\label{sec:discussion}
% ============================================================

% sections/ndss_07_discussion.tex  — NDSS 2027

\subsection{Limitations}

\paragraph{LSCPS as lower bound.}
Our rule-based capability extraction (\lscps) is a conservative
precision-recall tradeoff: it minimizes false positives at the cost
of missing obfuscated or dynamically-constructed dangerous calls.
Reported IC rates are therefore \emph{lower bounds} on true inconsistency.
CR = \MeanCR is an \emph{upper bound} on true disclosure rates.
LLM-based semantic adjudication (deferred to future work)
would expand detection coverage.

\paragraph{CAPE framework limitations.}
Of 85 pairs attempted at submission, 24 timed out under our pre-registered
175\,s hard-kill policy (LLM context compaction at high token usage).
We report both ITT (67.1\%) and completed-run (93.4\%) metrics side-by-side.
Controls and P2/P3 payload validations are 20-pair and 8-pair spot-checks;
expanding these to the full 125 pairs would further strengthen the causal claim.
We have not exhaustively tested all possible injection phrasings;
adversarial rephrasing likely increases the success rate further.

\paragraph{Single platform.}
Our measurement covers \clawhub and skills.sh only.
Other agent platforms (Cursor, Continue.dev, Claude Computer Use)
may have different transparency practices.
Cross-platform validation of skill.sh (90.9\% inconsistency)
suggests the finding generalizes.

\paragraph{Temporal snapshot.}
The ClawHub dataset reflects March 2026.
Skill descriptions may improve over time, though our temporal
analysis (\S\ref{sec:temporal}) shows the trend is toward
capability accretion without description update.

\subsection{Defenses}

\paragraph{Capability manifests.}
Marketplace platforms should require machine-readable capability
declarations (analogous to Android permissions) listing:
file-system access scope, network destinations, shell invocation flags,
credential access, and external API calls.
Our \lscps taxonomy provides a starting point for this schema.

\paragraph{Publish-time screening.}
Our \sysname pipeline can run at publish time to flag skills with
CR below a threshold (\eg CR $<$ 0.5), preventing the most egregious
mismatch patterns from reaching users.
At our measured FPR of $<$2\%, the cost is acceptable for platform operators.

\paragraph{Runtime capability monitoring.}
Agent runtimes should log and optionally confirm dangerous tool calls
(\texttt{exec}, \texttt{write}, \texttt{delete}) even when initiated by
installed skills, with user-visible notifications.
Our results show the agent's current safety model is heuristic-based
and inconsistent: shell exec is fully unguarded while obvious memory-write
requests are refused. A consistent capability-based policy is needed.

\paragraph{Trust-building attack countermeasure.}
The session warm-up bypass (\S\ref{sec:cape}) suggests that
agent safety policies should not rely on conversational trust context.
Capability whitelists established at skill-install time---not inferred
from session history---would prevent this bypass.

\subsection{Broader Implications}

The IC-1 finding (\ICone of dangerous skills omit capabilities) implies
that LLM agent marketplaces are currently operating on an \emph{honor system}
for capability disclosure.
As agent capabilities grow---computer use, multi-agent orchestration,
financial autonomy---the consequences of this trust gap compound.
Our work provides both the measurement evidence and the attack validation
needed to motivate platform-level policy change.


% ============================================================
\section{Related Work}
\label{sec:related}
% ============================================================

% sections/ndss_08_related.tex  — NDSS 2027

\paragraph{Prompt injection.}
Greshake \etal~\cite{greshake2023indirect} introduce indirect prompt injection
and demonstrate that LLM agents can be manipulated through
attacker-controlled content in retrieved documents.
Our work complements this: we show that when dangerous sink capabilities
are undisclosed, even a user who is aware of prompt injection
cannot protect themselves---they cannot know which installed skills
can be weaponized as sinks.
Recent work on multi-agent injection~\cite{wu2024autogen}
and tool-call hijacking~\cite{debenedetti2024agentdojo}
further establishes the attack surface; we provide the first
marketplace-scale measurement of the exposed capability surface.

Table~\ref{tab:injection_comparison} formally contrasts our threat model
with classical prompt injection and plugin-level attacks:

\begin{table}[h]
\centering
\caption{Source-sink composition attack vs.\ related threat models.}
\label{tab:injection_comparison}
\scriptsize
\setlength{\tabcolsep}{3pt}
\begin{tabular}{lllll}
\toprule
 & \textbf{Classic PI} & \textbf{Plugin atk} & \textbf{This work} \\
\midrule
Attacker controls & Input & Plugin code & Source skill output \\
Exploited gap & Model trust & Code execution & Capability omission \\
Requires runtime & Yes & Yes & No (install-time) \\
User can prevent? & Careful inputs & Vetted plugins & Not without disclosure \\
Sink needed? & No & No & Yes (IC-1 skill) \\
\bottomrule
\end{tabular}
\end{table}

\paragraph{Mobile app permission analysis.}
Extensive work on Android permissions~\cite{felt2011android,au2012pscout}
and privacy policy consistency~\cite{andow2019policylint,bui2021consistency}
motivated our methodology.
Liao \etal~\cite{liao2020measuring} measure over-privilege in Alexa skills
using a difference matrix approach; we adapt this to LLM agent marketplaces
where the permission model is replaced by natural-language descriptions.
Key differences: skills have richer implementation code (not just intents),
can execute arbitrary shell commands, and operate with full OS-level authority.

\paragraph{Software supply chain attacks.}
Ohm \etal~\cite{ohm2020backstabber} and Ladisa \etal~\cite{ladisa2023taxonomy}
taxonomy malicious packages in npm/PyPI.
LLM skill supply chain attacks differ in that the malicious payload
delivery mechanism is NLP-mediated: the LLM decides whether to call
a tool based on context, not static control flow.
This makes both detection (static analysis alone is insufficient)
and exploitation (prompt injection enables remote triggering) qualitatively different.

\paragraph{LLM agent security.}
Recent work studies jailbreaking~\cite{wei2023jailbroken},
tool misuse~\cite{ruan2024identifying},
and multi-agent trust~\cite{ju2024flooding}.
\sysname focuses on the marketplace supply chain rather than
the agent runtime---the attack surface created before any
user interaction, at skill install time.


% ============================================================
\section{Conclusion}
\label{sec:conclusion}
% ============================================================

% sections/ndss_09_conclusion.tex  — NDSS 2027

We presented \sysname, the first large-scale measurement of
capability-description inconsistency in LLM agent skill marketplaces,
and \cape, an empirical validation framework demonstrating
the real-world exploitability of the identified gaps.

Our analysis of \NskillsCH \clawhub skills reveals that \ICone
completely omit security-sensitive capabilities---shell execution,
credential access, file deletion, and data exfiltration---from visible
descriptions, yielding a mean Capability Recall of \MeanCR{}
(Cohen's $\kappa = \IAAsink$ inter-annotator agreement on ground truth).
This transparency failure is not latent risk.
\cape validates all \NpairsTotal source-sink pairs against real OpenClaw agents.
Under the pre-defined intent-to-treat metric (counting \NTimeout{} infrastructure
timeouts as failures), \NpairsSuccessITT/\NpairsTotal = \PairsSuccessRateITT{} succeed.
Under the completed-run metric (conditioned on agent returning a response),
\NpairsSuccessCR/\NpairsCompletedRun = \textbf{\PairsSuccessRateCR}
(95\% CI: \CICRlo--\CICRhi), with shell-execution sinks exploited at
\ShellSuccessRate{} (\NShellCompleted/\NShellCompleted{} completed pairs).
Only \NRefusal{} pairs (3.2\%) produce explicit agent refusals,
and these refusals are stochastic---the same task type succeeds in
67--100\% of other pairs---indicating the absence of a consistent safety policy.

Three structural findings stand out beyond the aggregate numbers.
First, the agent's safety policy is heuristic-based and inconsistently applied:
shell command execution via injected context is completely unguarded,
while direct memory-modification requests are occasionally refused
without a stable policy threshold.
Second, all three tested payload families (P1: \texttt{go vet}, P2: \texttt{node verify.js},
P3: \texttt{cat README}) produce comparable success rates,
confirming that the vulnerability is in the agent's trust model,
not in any property of a specific command syntax.
Third, the 64 infrastructure timeouts have a single root cause
(LLM context compaction at high prompt token usage)---
a measurement challenge, not an agent defense---and are being retried.

LLM skill marketplaces are operating on an honor system for
capability disclosure. As agent autonomy expands, this gap produces
compounding risk. We call on platform operators to adopt
machine-readable capability manifests, publish-time screening,
and consistent capability-based runtime policies as a minimum response.
Our tooling and dataset are released to facilitate this work.

\paragraph{Acknowledgments.}
Omitted for anonymous review.


\balance
\bibliographystyle{plain}
\bibliography{references}

\end{document}
